\chapter{第十章：决策的人机分工——什么时候用AI，什么时候用人}

# 第十章：决策的人机分工——什么时候用AI，什么时候用人

\textit{"指挥官的判断力是战争中最重要的因素。"}
——克劳塞维茨《战争论》

\hline\newpage
---

老王我最近腰疼，去医院拍了片子。放射科医生看了五分钟，说没事，开了点药让我回去了。

我心想，这五分钟电脑看了多久？不到三十秒。

那放射科医生五分钟在干嘛？他在跟AI"讨论"。AI先扫了一遍，标记了几个"可疑点"，医生再逐个过一遍，决定哪些是真的可疑、哪些是假警报。

这其实就是人机协作的一个缩影：AI做它擅长的——快速扫描、大海捞针；人做人擅长的——判断这个"针"到底是不是你要找的那根。

但问题是：什么时候该信AI，什么时候该信人？这个问题，比我腰疼还难治。

\section{一、一个医院的真实故事}

## 一、一个医院的真实故事

某家大型医院在2024年引入了一套AI辅助诊断系统。系统可以分析医学影像，帮助放射科医生识别可能的肿瘤、骨折和其他异常。

系统上线后，放射科医生很快发现了一个问题：系统的误报率比他们预期的要高。它标记了很多"可疑阴影"，但经过进一步检查发现大多数是良性。这让放射科医生陷入了困境：是相信自己的判断，还是相信系统的判断？

最终，他们采取了一个实用的方法：对于系统标记为"高风险"的案例，他们会特别仔细地再次检查。对于系统标记为"低风险"的案例，他们会相对快速地通过。这个方法在短期内有效——它让医生的工作效率提升了约30%，同时保持了诊断质量。

但这个案例揭示了一个更深层的问题：\textbf{人机协作的决策边界在哪里？}

不是所有的问题都适合由AI来回答，也不是所有的问题都需要人来判断。在某些情况下，AI比人强；在另一些情况下，人比AI强。但这种"谁更强"的判断本身，就需要人类的智慧。

克劳塞维茨说"指挥官的判断力是战争中最重要的因素"——在AI时代，这话照样成立，只不过"指挥官"得学会判断什么时候该信AI、什么时候该信自己。

\section{二、AI和人类，各自擅长什么}

## 二、AI和人类，各自擅长什么

要回答"什么时候用AI，什么时候用人"，首先需要理解AI和人类各自擅长什么。

这个问题的答案，不是简单粗暴的"AI擅长数据处理，人类擅长创意"能概括的。更准确的框架来自于对两种认知模式的区分：\textbf{模式识别}和\textbf{意义建构}。

\textbf{模式识别}是AI的强项。AI可以在海量数据中找到规律，可以识别图像、声音、文本中的特定模式，可以基于历史数据预测未来事件。这是AI在棋类游戏、医学影像诊断、语音识别、自然语言处理等领域超越人类的原因。

\textbf{意义建构}是人类的强项。人类能够理解上下文，能够建构意义，能够问出AI没有想过的问题，能够在不确定性中做出价值判断。这是AI在2024年仍然无法匹敌人类的领域。

但这个框架还不够精确。让我们更具体地分析几种典型的决策场景。

\textbf{场景一：已经有清晰标准的重复性决策}

贷款审批：申请人的收入、信用评分、负债比例——这些都是可量化的指标，批准或拒绝的标准是可以明确定义的。这类决策，AI做得比人更快、更准确、更少偏见。\textbf{答案：用AI。}

采用"能动式AI"的组织——让AI自主决策、持续迭代而非只是执行预设流程——效率提升中位数远高于传统自动化。这说明什么？当你把AI只当作"执行工具"，你只得到一小部分收益；当你把AI当作"决策参与者"，收益能翻番。当然，前提是你得有能力判断AI什么时候靠谱、什么时候不靠谱。

就像你ATM取钱，你希望机器给你数钱还是柜员给你数？机器又准又快，还不会跟你聊天耽误你时间。

\textbf{场景二：标准本身需要被建构的决策}

产品战略：下一个产品应该是什么？这不是一个有清晰输入和输出的问题。这是一个需要理解用户需求、判断市场趋势、整合多方利益的复杂问题。这个问题的"正确答案"在问题被解决之前是不知道的。\textbf{答案：用人，或者人机协作，但人的主导性更强。}

就像你问"中午吃什么"，这个问题有标准答案吗？没有。你得问自己想吃什么、今天心情怎么样、最近是不是在减肥、钱包里还有多少钱——这些信息综合起来，才形成一个决策。AI可以帮你查附近有什么馆子、评分怎么样、距离有多远，但"想吃什么"这件事，只有你自己知道。

\textbf{场景三：后果严重且不可逆的决策}

自动驾驶的事故决策：当车辆面临不可避免的碰撞时，应该优先保护车内乘客还是行人？这个问题没有数学上的最优解，因为它涉及价值判断。\textbf{答案：用人（或社会的法律和伦理框架），AI只能执行已由人类决定的规则。}

这就像电车难题——你没法让一个AI来做这个价值判断，因为它没有道德感，也没有法律主体资格。这个锅，AI背不了。

\textbf{场景四：需要跨领域整合的决策}

新市场进入：某个市场机会是否值得追求？这需要技术可行性评估、市场规模估算、竞争格局分析、组织能力匹配——这些是跨领域的综合判断。AI在单一领域可以提供很好的输入，但整合判断需要人类。\textbf{答案：人机协作，AI提供分析，人做出判断。}

就像你装修房子，AI能给你出水电图、结构图、效果图，但最后你得自己决定哪个方案住着舒服、哪个预算承受得起、哪个审美符合你的品味。

\section{三、决策的分级框架}

## 三、决策的分级框架

基于以上分析，我们可以建立一个决策分级框架。这个框架帮助组织系统性地思考哪些决策应该交给AI，哪些应该保留给人。

\textbf{第一级：AI全自动决策}

适用范围：输入可量化、标准可定义、错误可接受、反馈可快速获取的重复性决策。

典型案例：垃圾邮件识别、价格调整、库存补货推荐、内容推荐。

这一级决策，AI可以全自动完成，人类只需要在异常情况下介入。

就像你家的扫地机器人，它自己知道什么时候该扫、哪里该重点扫、什么时候该回去充电——你不需要盯着，你只管上班去，回来发现家干净了。

\textbf{第二级：AI建议+人类确认}

适用范围：输入复杂、需要上下文理解、后果较重、但标准相对清晰的决策。

典型案例：贷款审批（需要人工审核高风险案例）、医学影像诊断（需要放射科医生最终签字）、简历筛选（需要HR做最终判断）。

这一级决策，AI提供分析和建议，人类做出最终决定。人类的作用是审核AI的建议，识别AI可能忽略的异常情况。

就像你老婆让你下班顺便买菜，给你发了个购物清单。你到超市一看，"有机的"卖完了，你得自己决定买普通版的还是换个菜。这事AI清单搞不定，得你来。

\textbf{第三级：AI辅助+人类决策}

适用范围：输入高度复杂、涉及多方利益、后果严重、但可能没有明确最优解的决策。

典型案例：战略规划、新产品方向、高管任命、重大投资。

这一级决策，AI提供信息分析和选项建议，人类做出判断。关键是确保人类没有被AI的建议过度影响——AI可以扩展人类的视野，但不能替代人类的决策权。

就像你买房，AI能给你分析房价走势、周边配套、学区排名、通勤时间，但最后买不买、买哪套、多少钱买，这些决定还是你自己做。AI可以帮你看房，但它不能替你签字。

\textbf{第四级：人类决策，AI不参与}

适用范围：涉及根本性价值判断、涉及利益冲突、需要人类责任承担的决策。

典型案例：裁员决策、监管政策制定、伦理边界设定。

这一级决策，AI不应该介入或者只应该作为极其边缘的参考。因为这些决策需要人类的责任承担，AI不能被追究责任。

《论语》讲"不在其位，不谋其政"——AI也一样，不在其位，不担其责。你不能把裁员名单让AI来定，然后说出事了是AI的问题。锅就是你的，你推不掉。

\section{四、人机协作的常见陷阱}

## 四、人机协作的常见陷阱

在实践中，人机协作的决策往往比理论框架更复杂。以下是几个常见的陷阱。

\textbf{陷阱一：自动化偏见}

当AI系统给出建议时，人类倾向于盲目接受。这是"自动化偏见"——人们倾向于过度信任自动化系统的输出，即使在应该怀疑的时候也不怀疑。

这个陷阱的危险在于：当AI出错时，盲目信任AI的人类会把AI的错误放大，而不是纠正它。在高风险决策场景中，这可能是灾难性的。

避免这个陷阱，需要建立一种"批判性信任"的文化——对AI的建议保持警觉，但不盲目拒绝；接受建议，但同时独立思考。

老话说"尽信书不如无书"——尽信AI不如无AI。你得把AI当一个特别聪明的同事看，而不是当先知看。同事可能出错，先知也可能出错，你得自己动脑子。

\textbf{陷阱二：决策责任模糊}

当AI参与了决策过程时，谁对决策负责？如果AI的建议导致了负面结果，是AI的问题？是开发AI的团队的问题？还是采纳建议的人的问题？

这个问题在法律和伦理上都还没有清晰的答案。但在组织层面，你需要尽可能地提前明确这个问题的答案。

实践中的建议是：无论AI参与了多少，最终决策的人类永远承担最终责任。AI是工具，工具的制造者和使用者共同对使用工具的决策负责。

你雇人干活，出了问题你负责，不是工人负责——这个道理在AI时代也一样。你用了AI的建议，采纳了，出了事就是你负责。你不能说"是AI说的"，然后把自己摘干净。

\textbf{陷阱三：AI能力的过度信任和过度不信任}

过度信任AI：认为AI永远是对的，忽视AI的局限性。

过度不信任AI：认为AI永远不可靠，坚持所有事情都由人来做。

两种极端都是有害的。过度信任会导致错误决策，过度不信任会导致效率损失和组织能力的停滞。

健康的做法是：基于具体场景和AI在特定任务上的历史表现，来判断应该给予多大程度的信任。

就像你骑自行车，你不会因为有一次摔倒了就永远不骑车，你也不会因为骑车很方便就闭着眼睛骑。根据路况、天气、你的状态来决定骑多快、什么时候下来推着走，这才健康。

\section{五、决策质量评估的闭环}

## 五、决策质量评估的闭环

人机协作决策的有效性，需要通过持续评估来保证。

传统的决策评估是滞后的——一个决策做出后，需要等几周、几个月甚至几年才能知道结果。当AI介入决策时，这个滞后性让评估变得更加困难——你需要知道在没有AI的情况下，同样的决策会是怎么样的，才能判断AI是否真的有帮助。

建立决策质量评估的闭环，需要：

\textbf{第一，建立可追溯的决策记录。} 每一次人机协作的决策，都需要记录AI的建议是什么、人类的决定是什么、最终的决策结果是什么。这些记录是后续分析的基础。

\textbf{第二，定期审计高风险决策。} 对于后果严重的决策，需要有定期的审计机制，检查AI的建议和人类的决定之间是否存在系统性的偏差，是否存在人类过度依赖或过度忽视AI建议的模式。

\textbf{第三，建立反馈机制。} 当决策结果出来后，需要有机制让这个结果反馈到AI系统的改进中。如果AI在某个类型的决策上持续表现不佳，需要有系统性的方法来识别和修复这个问题。

就像你开车，不是上了路就完事了，你得看导航、听语音提示、根据自己的判断决定要不要按AI说的走、最后还得复盘这条路走得对不对。人机协作也一样，是个持续反馈的闭环。

\section{六、结语：人机分工是持续演化的动态边界}

## 六、结语：人机分工是持续演化的动态边界

关于"什么时候用AI，什么时候用人"这个问题，没有一劳永逸的答案。

因为AI的能力在持续进化。2023年认为只有人类能做的事情，2025年可能已经有AI可以做得很好。即使在同一类任务上，AI的能力边界也在持续移动。

这意味着，人机分工的边界，需要持续被重新评估。去年的最佳实践，可能今年就过时了。

管理者需要建立的，不是一套静态的"AI做这个，人类做那个"的规则，而是一套持续评估和调整人机分工的机制。

这套机制需要包括：
\begin{itemize}
\item 持续追踪AI能力边界的演变
\item 定期评估当前的人机分工是否仍然是最优的
\item 建立实验和学习文化，持续测试人机协作的新方式

\end{itemize}
人机分工，不是关于谁更强的问题，而是关于如何最优组合的问题。在不同的场景下，最优组合是不同的。而最优组合本身，也在随着技术和环境的变化而持续变化。

这是AI时代管理者面临的一个全新的课题。

《孙子兵法》讲"兵无常势，水无常形"——人机分工也没有一成不变的公式。你得根据敌情（场景）来调整部署（分工），根据战局（技术发展）来改变策略（边界）。这不是一次性的事，是天天都得想的事。

\hline\newpage
---

\textbf{本章小结：}

1. AI擅长模式识别，人类擅长意义建构——但这个区分在不同场景下需要具体分析
2. 决策分级框架：AI全自动决策 → AI建议+人类确认 → AI辅助+人类决策 → 人类决策AI不参与
3. 人机协作的三个陷阱：自动化偏见（盲目信任AI）、决策责任模糊、对AI能力的过度信任或过度不信任
4. 建立决策质量评估闭环：可追溯的决策记录、高风险决策定期审计、决策结果的反馈机制
5. 人机分工是动态边界——AI能力在持续进化，分工边界需要持续重新评估，不能一劳永逸
